\documentclass[11pt]{article}

\usepackage[a4paper,margin=0.78in]{geometry}
\usepackage[T1]{fontenc}
\usepackage[utf8]{inputenc}
\usepackage{lmodern}
\usepackage{amsmath,amssymb,amsthm,mathtools}
\usepackage{enumitem,booktabs,array,tabularx,microtype,xcolor}
\usepackage[numbers,sort&compress]{natbib}
\usepackage[colorlinks=true,linkcolor=blue!65!black,
  citecolor=blue!65!black,urlcolor=blue!70!black]{hyperref}
\usepackage[nameinlink,noabbrev]{cleveref}
\setlength{\emergencystretch}{3em}

\hypersetup{
  pdftitle={Multiplicative Quotient Lifts on Determinant-Two Fibers},
  pdfauthor={Liang Wang},
  pdfsubject={Exact Mobius recovery and stable nonpretentiousness},
  pdfkeywords={Mobius function, multiplicative functions, determinant two, pretentious distance}
}

\newtheorem{theorem}{Theorem}[section]
\newtheorem{proposition}[theorem]{Proposition}
\newtheorem{corollary}[theorem]{Corollary}
\newtheorem{lemma}[theorem]{Lemma}
\theoremstyle{definition}
\newtheorem{definition}[theorem]{Definition}
\theoremstyle{remark}
\newtheorem{remark}[theorem]{Remark}

\newcommand{\Lzero}{\mathrm{L0}}
\newcommand{\Lone}{\mathrm{L1}}
\newcommand{\Ltwo}{\mathrm{L2}}
\newcommand{\N}{\mathbb N}
\newcommand{\PROVED}{\textnormal{\textsc{proved}}}
\newcommand{\OPEN}{\textnormal{\textsc{open}}}
\newcolumntype{Y}{>{\raggedright\arraybackslash}X}

\title{\textbf{Multiplicative Quotient Lifts on}\\
\textbf{Determinant-Two Fibers: Exact M\"obius}\\
\textbf{Recovery and Stable Nonpretentiousness}}
\author{Liang Wang\\
School of Artificial Intelligence and Automation,\\
Huazhong University of Science and Technology, Wuhan 430074, P.R. China\\
\texttt{wangliang.f@gmail.com}}
\date{July 2026}

\begin{document}
\maketitle

\begin{abstract}
The determinant-two pullback previously represented
\(\mu(D)\mu(V)\) as a Liouville shift-two product multiplied by two
quotient-squarefree masks.  We give a different exact
representation.  For every positive integer \(c\), an explicit
one-bounded multiplicative function \(G_c\) satisfies
\[
 G_c(cm)=\mu(m).
\]
If \(D(z)=d+sz\), \(V(z)=u+az\), and \(su-ad=2\), then with
\(t=aD(z)=ad+asz\),
\[
 \mu(D(z))\mu(V(z))=G_a(t)G_s(t+2).
\]
Thus the full squarefree information is already inside two
multiplicative functions; no prime-square cutoff or CRT squarefree
tail is needed for this two-point core.

At prime arguments, \(G_c\) differs from the Liouville function only
at primes \(p\) with \(p\parallel c\).  Its pretentious distance is
therefore bounded below by the Liouville distance minus
\(2\sum_{p\parallel c}p^{-1}\).  For small-polylogarithmic \(c\) this
loss is lower order than the known logarithmic lower bound.  The
construction supplies an exact \(\Lzero\) lift and a uniform
\(\Lone\) nonpretentious input.  It does not supply physical weights,
generic phases, all prefixes, a four-point estimate, positive
\(\Ltwo\), or a prime-pair theorem.
\end{abstract}

\section{The quotient lift}

For a prime \(p\), write \(e_p=v_p(c)\).  We define the local values
of \(G_c\) before taking their multiplicative product.

\begin{definition}[Quotient-M\"obius lift]
\label{def:Gc}
For \(c\ge1\), let \(G_c:\N\to\{-1,0,1\}\) be the multiplicative
function determined by
\begin{equation}
 G_c(p^j)=
 \begin{cases}
  (-1)^j,&0\le j<e_p,\\
  \mu(p^{j-e_p}),&j\ge e_p.
 \end{cases}
\label{eq:local}
\end{equation}
\end{definition}

The two cases agree with \(G_c(1)=1\).  This is a multiplicative
definition through independent prime-power factors.  It is generally
not completely multiplicative.

\begin{theorem}[Exact quotient recovery]
\label{thm:quotient}
For every \(c,m\ge1\),
\begin{equation}
 \boxed{G_c(cm)=\mu(m).}
\label{eq:quotient}
\end{equation}
\end{theorem}

\begin{proof}
For each prime \(p\), the exponent of \(p\) in \(cm\) is
\(e_p+v_p(m)\), which lies in the second case of
\eqref{eq:local}.  Its local factor is
\(\mu(p^{v_p(m)})\).  Multiplication over primes gives \(\mu(m)\).
\end{proof}

\begin{remark}[Why the first case is retained]
Only values on multiples of \(c\) are needed for
\eqref{eq:quotient}; the values below \(e_p\) may therefore be
chosen.  The choice \((-1)^j\) makes the prime values agree with
\(\lambda(p)=-1\) unless \(e_p=1\), which minimizes the
pretentious-distance perturbation used below.
\end{remark}

\section{The determinant-two fiber}

Let
\begin{equation}
 D(z)=d+sz,\qquad V(z)=u+az,\qquad su-ad=2,
\label{eq:fiber}
\end{equation}
with positive integral values on the interval under consideration.
In the TPC carrier, \(a,s\) are positive, odd and coprime
\citep{WangTPC127}.  Put
\begin{equation}
 Q_{\rm fib}=as,\qquad t(z)=aD(z)=ad+Q_{\rm fib}z.
\label{eq:t}
\end{equation}

\begin{theorem}[Exact determinant-two M\"obius lift]
\label{thm:fiber}
For every integer \(z\) on which the forms are positive,
\begin{align}
 t(z)+2&=sV(z),\label{eq:shift}\\
 t(z)&\equiv ad\pmod{Q_{\rm fib}},\label{eq:progression}\\
 \mu(D(z))\mu(V(z))
 &=\boxed{G_a(t(z))G_s(t(z)+2)}
\label{eq:fiber-lift}
\end{align}
\end{theorem}

\begin{proof}
The determinant equation gives
\[
 sV(z)-aD(z)=su-ad=2,
\]
which proves \eqref{eq:shift}; \eqref{eq:progression} follows from
\eqref{eq:t}.  Both \(a\mid t(z)\) and \(s\mid t(z)+2\).
\Cref{thm:quotient} gives
\[
 G_a(t(z))=\mu(t(z)/a)=\mu(D(z)),
 \quad
 G_s(t(z)+2)=\mu((t(z)+2)/s)=\mu(V(z)).
\]
\end{proof}

\begin{corollary}[No squarefree truncation in the two-point core]
\label{cor:no-tail}
The right side of \eqref{eq:fiber-lift} contains the full M\"obius
values.  Applying a two-point theorem directly to \(G_a(t)G_s(t+2)\)
requires neither quotient-squarefree masks nor a prime-square cutoff.
Consequently it creates no squarefree truncation tail.
\end{corollary}

This is a new representation of the same literal core; it does not
delete the remaining coprimality, periodic, physical or prefix
factors of the full packet.

\section{Pretentious distance under the lift}

For one-bounded multiplicative functions, write
\[
 M(g;Y,Q)
 =
 \inf_{\substack{|v|\le Y\\q\le Q\\\chi\ ({\rm mod}\ q)}}
 \sum_{p\le Y}
 \frac{1-\operatorname{Re}(g(p)\overline{\chi(p)}p^{-iv})}{p}.
\]
The convention is the squared pretentious distance used in
\citet{TaoTeravainen2026}.

\begin{lemma}[Prime-level perturbation]
\label{lem:prime}
For every prime \(p\),
\[
 G_c(p)=
 \begin{cases}
  1,&p\parallel c,\\
  -1,&p\nparallel c.
 \end{cases}
\]
Thus \(G_c(p)\ne\lambda(p)\) precisely when \(p\parallel c\).
\end{lemma}

\begin{proof}
If \(v_p(c)=0\), the second line of \eqref{eq:local} gives
\(\mu(p)=-1\).  If \(v_p(c)=1\), it gives \(\mu(1)=1\).
If \(v_p(c)\ge2\), the first line gives \((-1)^1=-1\).
\end{proof}

\begin{proposition}[Uniform nonpretentious stability]
\label{prop:distance}
For all \(Y,Q\ge2\),
\begin{equation}
 \boxed{
 M(G_c;Y,Q)
 \ge M(\lambda;Y,Q)
 -2\sum_{\substack{p\le Y\\p\parallel c}}\frac1p.}
\label{eq:distance}
\end{equation}
\end{proposition}

\begin{proof}
Fix one \(\chi\) and \(v\) in the defining infimum.  By
\cref{lem:prime}, the two prime sums agree outside \(p\parallel c\).
At an exceptional prime, changing one unit-disc value to another
changes the real part by at most two.  Hence the \(G_c\) sum is at
least the \(\lambda\) sum minus the displayed error.  Taking the
infimum proves \eqref{eq:distance}.
\end{proof}

\begin{lemma}[Small-polylogarithmic perturbations are lower order]
\label{lem:reciprocal}
For \(c\ge3\),
\[
 \sum_{p\mid c}\frac1p\ll1+\log_3^+ c,
 \qquad
 \log_3^+c:=\max\{0,\log\log\log c\}.
\]
In particular, if \(c\le(\log X)^A\) for fixed \(A\), then this sum
is \(o(\log_2X)\).
\end{lemma}

\begin{proof}
For fixed size of the product of the distinct prime divisors, the
reciprocal sum is maximized by the smallest primes.  The standard
primorial and Mertens estimates then give
\(\sum_{p\mid c}p^{-1}\ll1+\log_3^+c\).  The stated consequence is
immediate.
\end{proof}

\Citet[equation~(1.12)]{MatomakiRadziwillTao2015} prove, uniformly
for characters of modulus at most a fixed small power of \(\log X\),
\begin{equation}
 M(\lambda;X^2,\log^{1/125}X)
 \ge\left(\frac13-o(1)\right)\log_2X.
\label{eq:lambda-distance}
\end{equation}
Combining the preceding results gives the needed source condition.

\begin{corollary}[Callable nonpretentious branch]
\label{cor:callable}
There is an absolute \(\alpha>0\) such that, for every fixed
\(A>0\) and all sufficiently large \(X\) depending on \(A\),
uniformly for \(c\le(\log X)^A\),
\[
 \exp M(G_c;X^2,\log^{1/125}X)
 \gg(\log X)^\alpha.
\]
After decreasing \(\alpha\), the nonpretentious hypothesis of
\citet[Theorem~3.1]{TaoTeravainen2026} holds with
\(\mathcal L=(\log X)^\alpha\).
\end{corollary}

\begin{proof}
Equations \eqref{eq:distance}--\eqref{eq:lambda-distance} and
\cref{lem:reciprocal} give
\[
 M(G_c;X^2,\log^{1/125}X)
 \ge(1/3-o(1))\log_2X.
\]
Exponentiate and choose any fixed \(\alpha<1/3\), with room for the
implicit source constant.
\end{proof}

\section{Scope and machine certificate}

The deterministic script evaluates \eqref{eq:local} exactly,
verifies \eqref{eq:quotient} for thousands of integer pairs, checks
ordinary multiplicativity on coprime inputs, and verifies
\eqref{eq:fiber-lift} on several determinant-two fibers.  It records
the prime modification set \(\{p:p\parallel c\}\).  These finite
checks are regressions for the algebra; they do not numerically prove
\eqref{eq:lambda-distance}.

\begin{center}
\small
\begin{tabularx}{\textwidth}{@{}p{0.48\textwidth}Y@{}}
\toprule
Statement & Level and status\\
\midrule
\(G_c(cm)=\mu(m)\) & Exact \(\Lzero\), \(\PROVED\).\\
Determinant-two identity \eqref{eq:fiber-lift} &
Exact \(\Lzero\), \(\PROVED\).\\
Pretentious stability \eqref{eq:distance} &
Unconditional comparison, \(\PROVED\).\\
Small-polylog source condition &
Sourced quantitative \(\Lone\), \(\PROVED\).\\
Bounded periodic reassembly &
Available through TPC-147 \citep{WangTPC147}.\\
Physical multiplier, generic phase, all prefixes &
\(\OPEN\); not supplied.\\
Four-point input or positive \(\Ltwo\) &
Not proved.\\
Fixed \(X\)-power, \(1/400\), prime pairs &
Not proved.\\
\bottomrule
\end{tabularx}
\end{center}

The lift therefore removes the squarefree-CRT gate for the two-point
M\"obius core without claiming that a two-point good-scale theorem
is the complete H3 input.

\begingroup
\footnotesize
\raggedright
\bibliographystyle{plainnat}
\bibliography{references}
\endgroup

\end{document}
